\section{Stage 3: Results}
\label{sec:results}

\IfFileExists{generated/main_table.tex}{\begin{table*}[t]
\centering
\caption{Reproduced test-set results on APTOS 2019 (550 test images, faithful split). Paper: QWK 0.954 (single model, ImageNet init) and 0.967 (EyePACS init + snapshot ensemble).}
\label{tab:main}
\footnotesize
\begin{tabular}{llcccccc}
\toprule
Configuration & Prediction & QWK & Acc. & Macro-F1 & Macro AUC & Minority recall & $|\Delta|\geq2$ errors \\
\midrule
300\,px/60 ep, lr $10^{-3}$ (paper) & final epoch & 0.808 & 0.785 & 0.536 & 0.887 & 0.159 & 37 \\
 & best val QWK & 0.858 & 0.780 & 0.512 & 0.914 & 0.211 & 31 \\
 & snapshot vote & 0.846 & 0.796 & 0.527 & 0.923 & 0.205 & 37 \\
 & snapshot mean & 0.839 & 0.796 & 0.527 & 0.923 & 0.193 & 38 \\
\midrule
300\,px/60 ep, lr $10^{-4}$ & final epoch & 0.858 & 0.809 & 0.610 & 0.899 & 0.388 & 40 \\
 & best val QWK & 0.860 & 0.813 & 0.599 & 0.908 & 0.342 & 27 \\
 & snapshot vote & 0.879 & 0.822 & 0.623 & 0.934 & 0.359 & 29 \\
 & snapshot mean & 0.882 & 0.825 & 0.629 & 0.934 & 0.370 & 28 \\
\midrule
224\,px/30 ep, lr $10^{-3}$ (M2) & final epoch & 0.827 & 0.771 & 0.473 & 0.897 & 0.148 & 48 \\
 & best val QWK & 0.827 & 0.771 & 0.473 & 0.897 & 0.148 & 48 \\
 & snapshot vote & 0.788 & 0.767 & 0.412 & 0.913 & 0.000 & 50 \\
 & snapshot mean & 0.791 & 0.769 & 0.416 & 0.913 & 0.000 & 49 \\
\midrule
224\,px/30 ep, lr $10^{-4}$ (M2) & final epoch & 0.862 & 0.798 & 0.603 & 0.897 & 0.370 & 30 \\
 & best val QWK & 0.887 & 0.815 & 0.630 & 0.894 & 0.428 & 24 \\
 & snapshot vote & 0.887 & 0.822 & 0.613 & 0.930 & 0.347 & 27 \\
 & snapshot mean & 0.887 & 0.820 & 0.611 & 0.930 & 0.347 & 27 \\
\bottomrule
\end{tabular}
\end{table*}
}{}
\IfFileExists{generated/per_class_table.tex}{\begin{table}[t]
\centering
\caption{Per-class test metrics, exact paper protocol (300\,px, 60 epochs, lr $10^{-3}$).}
\label{tab:perclass}
\footnotesize
\begin{tabular}{llcccc}
\toprule
Grade & Model & Prec. & Recall & F1 & $n$ \\
\midrule
No DR & best val & 0.97 & 0.96 & 0.96 & 271 \\
Mild & best val & 0.44 & 0.21 & 0.29 & 56 \\
Moderate & best val & 0.65 & 0.93 & 0.76 & 150 \\
Severe & best val & 0.27 & 0.10 & 0.15 & 29 \\
PDR & best val & 0.52 & 0.32 & 0.39 & 44 \\
\midrule
No DR & snapshot vote & 0.94 & 0.97 & 0.96 & 271 \\
Mild & snapshot vote & 0.50 & 0.38 & 0.43 & 56 \\
Moderate & snapshot vote & 0.70 & 0.91 & 0.79 & 150 \\
Severe & snapshot vote & 0.00 & 0.00 & 0.00 & 29 \\
PDR & snapshot vote & 0.53 & 0.41 & 0.46 & 44 \\
\bottomrule
\end{tabular}
\end{table}
}{}

\subsection{Faithful reproduction (exact protocol)}
Table~\ref{tab:main} summarises the test results. With the paper's protocol and hyper-parameters, the
best-validation model (epoch 15) reaches QWK~\FullBestQWK{} (accuracy \FullBestAcc{}, macro-F1 \FullBestMF{}),
against 0.954 reported in~\cite{chilukoti2024}: a gap of about 0.10 QWK. The final-epoch model is weaker
(QWK \FullFinalQWK{}). The 10-snapshot majority vote, the paper's central contribution, gives \FullVoteQWK{}:
better than the final epoch but below the best single checkpoint.

Training was unstable (Fig.~\ref{fig:curves}, Fig.~\ref{fig:compare}): validation QWK fell below 0.75 at
epochs 3, 5, 8, 10, 17, 22, 23 and 36 (to 0.556 at epoch 36) and recovered within one or two epochs, and the
best validation score (0.895 at epoch 15) was never reached again. Test QWK of the individual snapshots
(epochs 6, 12, \dots, 60) was 0.830, 0.770, 0.828, 0.834, 0.803, 0.522, 0.808, 0.773, 0.781 and 0.808. The
members are of uneven quality and one is very poor, so the vote cannot exceed the best member. Snapshot
ensembling assumes comparably good, diverse members~\cite{huang2017snapshot}, which an unstable trajectory
without a cyclic learning-rate schedule does not provide.

\begin{figure}[t]
\centering
\includegraphics[width=\columnwidth]{s2_paper300_curves.pdf}
\caption{Training loss and validation metrics, exact protocol with the paper's lr $10^{-3}$.}
\label{fig:curves}
\end{figure}

\subsection{Per-grade behaviour}
QWK and accuracy hide a severe minority-grade failure (Table~\ref{tab:perclass}, Fig.~\ref{fig:cm}).
No DR (recall 0.96) and Moderate (0.93) are learned well, but \textbf{Severe NPDR is almost never recognised}:
3 of 29 Severe eyes for the single model, 0 for the vote, and at most 5 in any snapshot. PDR recall is
\FullBestPDRRec{} (single model) and \FullVotePDRRec{} (vote), and Mild recall only \FullBestMildRec{}.
Severe cases are mostly called Moderate (17/29) or PDR (9/29), 24 of 44 PDR eyes are called Moderate, and
34 of 56 Mild eyes are called Moderate. The errors are mostly adjacent-grade, which is why QWK stays near
0.86 even though two of five grades are essentially unrecognised; this is exactly the ``accuracy (and even QWK)
can be misleading'' gap identified in our literature review.

\begin{figure}[t]
\centering
\includegraphics[width=0.78\columnwidth]{s2_paper300_best_val_cm.pdf}
\caption{Test confusion matrix of the reproduced single model (exact protocol, paper hyper-parameters).}
\label{fig:cm}
\end{figure}

\IfFileExists{sections/diagnostic.tex}{\subsection{Diagnostic: learning rate $10^{-4}$}
\label{sec:diag}
To test whether optimisation explains the gap, we repeated the exact protocol with a single change, lr
$10^{-3}\to10^{-4}$ (same split, seed, resolution, epochs and snapshots). Training became stable: validation QWK
stayed between 0.81 and 0.89 for all 60 epochs with no collapses (Fig.~\ref{fig:compare}), and training loss
fell to 0.12 by epoch 10 (0.55 at lr $10^{-3}$). The best-validation checkpoints are similar (QWK
\FulllowBestQWK{} vs.\ \FullBestQWK{}), but every other prediction improves: final epoch \FullFinalQWK{}$\to$\FulllowFinalQWK{},
snapshot vote \FullVoteQWK{}$\to$\FulllowVoteQWK{} and snapshot mean \FullMeanQWK{}$\to$\FulllowMeanQWK{}
(macro-F1 \FulllowMeanMF{}). Minority recall rises from \FullBestMinRec{} to \FulllowMeanMinRec{} (ensemble mean;
Severe \FulllowMeanSevRec{}, PDR \FulllowMeanPDRRec{}).

With stable training the snapshots are uniformly good (test QWK 0.823--0.885) and the ensemble now beats the
best single model, consistent with the snapshot-ensemble assumption that members are of comparable quality.
Severe NPDR remains the weakest grade (5 of 29 recognised by the ensemble; 16 called Moderate) and Mild is still
mostly confused with Moderate (26/56). The low training loss with a flat validation curve suggests the model
overfits after about 10 epochs, which a learning-rate schedule or stronger augmentation could address.
This run is the baseline we carry into Stage~4.

\subsection{Effect of resolution and schedule}
\label{sec:res}
Comparing the exact protocol with our compute-limited pilot isolates the effect of using 224\,px for 30 epochs
(Table~\ref{tab:main}, Fig.~\ref{fig:compare}). At lr $10^{-3}$ the full protocol helps: best single model
\SIIBestQWK{}$\to$\FullBestQWK{} and vote \SIIVoteQWK{}$\to$\FullVoteQWK{}, mostly because the vote is no longer
dragged down as far by collapsed snapshots. At lr $10^{-4}$ it does not: the pilot's best model and vote
(\SIIlowBestQWK{}, \SIIlowVoteQWK{}) are at least as good as the exact protocol's (\FulllowBestQWK{},
\FulllowVoteQWK{}). Resolution and schedule therefore account for at most $\approx$0.03 QWK, and the learning
rate matters more than either.

\begin{figure}[t]
\centering
\includegraphics[width=\columnwidth]{val_qwk_compare.pdf}
\caption{Validation QWK per epoch for all four runs. lr $10^{-3}$ (red) is unstable at both resolutions;
lr $10^{-4}$ (blue) is stable. The dashed line is the paper's reported test QWK.}
\label{fig:compare}
\end{figure}
}{}

\subsection{Reproduction-gap analysis}
\label{sec:gap}
\begin{itemize}
\item \textbf{Duplicate leakage does not explain it.} Excluding the 29 test images with an identical copy in
training changes the single model's QWK from 0.858 to 0.862 (0.879 to 0.879 for the lr-$10^{-4}$ vote). The
leaked images are actually predicted \emph{worse} (58.6\% accuracy), because several copies carry conflicting grades.
\item \textbf{Optimisation (learning rate) explains part of it.} Full fine-tuning with Adam at $10^{-3}$ and
gradient clipping at 0.1 is unstable at both resolutions and batch sizes; lowering it to $10^{-4}$ removes the
collapses and improves the final-epoch and ensemble predictions by 0.03--0.05 QWK (Sec.~\ref{sec:diag}). The paper's own Table~3, where VGG and ResNet models
predict only class~0 under the same settings, is consistent with this.
\item \textbf{Resolution and schedule explain little.} Moving from our 224\,px / 30-epoch pilot to the exact
300\,px / 60-epoch protocol adds 0.03 QWK at lr $10^{-3}$ and nothing at $10^{-4}$ (Sec.~\ref{sec:res}).
\item \textbf{Unreported details.} The split, loss, CLAHE parameters, and whether the reported model is the
final or the best-validation one are not stated; each can move QWK by a few points on a 550-image test set.
\item \textbf{Reported numbers.} The internal inconsistencies noted in Sec.~2 (e.g.\ identical precision, recall
and QWK; macro AUC below 0.5) mean the reported 0.954 may itself not be reliable. With the stated protocol run
exactly, a gap of $\approx$0.07 QWK remains that none of the stated settings account for.
\end{itemize}

\subsection{Failure cases}
Fig.~\ref{fig:fail} shows the most confidently wrong test predictions of the exact-protocol model. Eight of the
ten are PDR eyes predicted as Mild or Moderate with high confidence, including one with a large pre-retinal
haemorrhage; the other two are a dim, low-contrast No DR image and a Mild image both predicted as PDR.
Neovascularisation is small and easy to miss, and the model appears to rely on exudates and overall image
appearance instead.

\begin{figure*}[t]
\centering
\includegraphics[width=0.9\textwidth]{s2_paper300_failures.pdf}
\caption{Ten largest-error test predictions of the reproduced single model, exact protocol (true grade,
predicted grade, confidence).}
\label{fig:fail}
\end{figure*}

\section{Limitations}
Single dataset (APTOS~2019) and a single seed per configuration; no external validation; the EyePACS-pretrained
configuration (0.967) was not reproduced. The 550-image test set contains only 29 Severe and 44 PDR images, so
per-grade recall estimates are noisy, and QWK differences of 0.01--0.02 between configurations are within
run-to-run variation. This is a reproduction study; no clinical claims are made.

\section{Next Step (Stage 4 preview)}
The reproduction shows that the paper's pipeline fails on exactly the grades that matter most clinically. Our
Stage~4 hypothesis will therefore change one variable, the training loss, from cross-entropy to a
class-balanced, ordinal-aware objective, on top of the stable lr-$10^{-4}$ exact-protocol baseline (snapshot
ensemble QWK \FulllowMeanQWK{}, Severe recall \FulllowMeanSevRec{}), and will be stated with a predicted
direction and magnitude before any code is run.
